\color{unidarkgreen}{\subsection{Контроль силового выключателя}}
\color{black}

\begin{enumerate}[label=\arabic{section}.\arabic{subsection}.\arabic*,  align=left, labelsep=4pt, leftmargin=0pt, itemindent=57pt]

\item
Функциональный блок контроля силового выключателя (КСВ) предназначен для обработки
технологической информации от выключателя и его цепей управления, а также от трансформатора тока,
позволяя определять их состояние, информировать о возможных неисправностях и формировать
соответствующую сигнализацию.

\begin{figure}[H]
\centering
\includegraphics[width=0.8\textwidth,height=0.8\textheight,keepaspectratio]{\fbpath/761.3371 КСВ/img/CTRL_Q.pdf}
\captionof{figure}{Функциональная схема контроля исправности выключателя}\label{pic_ksv:CTRL_Q}
\end{figure}

\item
Контроль готовности привода выключателя осуществляется по внешнему сигналу \textit{«Пруж. незаведены
В»}. Сигнал о незаведенном состоянии пружины \textit{«Ср. Пруж. не заведены В»} формируется с выдержкой времени
\textbf{«Тср Пруж. не заведены В»}. 

\item
Контроль целостности цепей управления осуществляется по сигналам состояния цепей первой и второй
группы электромагнитов отключения (ЭМО1 и ЭМО2) и цепи электромагнита включения (ЭМВ). Наличие второй
группы электромагнитов отключения определяется положением программного переключателя
\textbf{«Кол-во ЭМО В»}. При одновременном наличии или отсутствии сигналов контроля цепей ЭМО и ЭМВ в течение
регулируемой выдержки времени \textbf{«Тср Неиспр. ЦУ В»} формируется сигнал \textbf{«Неиспр. ЦУ В»}. Выдержка
времени должна быть не менее времени взвода пружины привода выключателя.

\item
Обогрев выключателя контролируется по внешнему сигналу \textit{«Неиспр. обогр. В»}. Соответствующий сигнал
о неисправности обогрева \textit{«Ср. Неиспр. обогр. В»} формируется с регулируемой выдержкой времени \textbf{«Тср
Неиспр. обогр. В»}.

\item
Сигнализация о неисправности выключателя (сигнал \textit{«Неиспр. В»}) срабатывает при наличии одного из
сигналов: \textit{«Неиспр. ЦУ В»}, \textit{«Ср. Пруж. не заведены В»}, \textit{«Ср. Низк. давл. SF6 В»}, \textit{«Ср. Авар. давл. SF6 В»}, \textit{«Ср.
Автомат ШП В»}, \textit{«Ср. Неиспр. обогр. В»}, а также от сигнала \textit{«Ср. ЗНФ»}.

\begin{figure}[H]
\centering
\includegraphics[width=0.8\textwidth,height=0.8\textheight,keepaspectratio]{\fbpath/761.3371 КСВ/img/CTRL_CT_Q.pdf}
\captionof{figure}{Функциональная схема контроля исправности ТТ}\label{pic_ksv:CTRL_CT_Q}
\end{figure}

\item
В устройстве также предусмотрен контроль состояния ТТ. Контроль давления элегаза в ТТ
осуществляется по внешним сигналам \textit{«Низк. давл. SF6 ТТ»} и \textit{«Авар. давл. SF6 ТТ»}. Соответствующие сигналы
о низком и аварийном давлении элегаза \textit{«Ср. Низк. давл. SF6»} и \textit{«Ср. Авар. давл. SF6»} формируются с
выдержками времени, равными 1 с. Обогрев ТТ контролируется по внешнему сигналу
\textit{«Неиспр. обогр. ТТ»}. Соответствующий сигнал о неисправности обогрева \textit{«Ср. Неиспр. обогр. ТТ»}
формируется с регулируемой выдержкой времени \textbf{«Тср Неиспр. обогр. ТТ»}. По наличию одного из указанных
сигналов формируется сигнал о неисправности ТТ \textit{«Неиспр. ТТ»}. При вводе программного ключа \textbf{«Откл. от
Авар. давл. SF6 ТТ»} предусматривается действие на отключение выключателя при аварийном снижении
давления элегаза в ТТ.

\begin{figure}[H]
\centering
\includegraphics[width=0.8\textwidth,height=0.8\textheight,keepaspectratio]{\fbpath/761.3371 КСВ/img/CTRL_OC_Q.pdf}
\captionof{figure}{Функциональная схема контроля оперативного тока цепей сигнализации}\label{pic_ksv:CTRL_OC_Q}
\end{figure}

\item
Контроль отсутствия питания цепей сигнализации выключателя и ТТ осуществляется по внешним сигналам \textit{«Контр.
опертока ЭМВ/ЭМО1 В»}, \textit{«Контр. опертока ЭМО2 В»} и \textit{«Контр. опертока сигн. В»}. При состоянии «0» одного
из сигналов срабатывает сигнализация о неисправности цепей сигнализации выключателя и ТТ (сигнал
\textit{«Неиспр. сигн. В/ТТ»}) с выдержкой времени 1 с.  

\begin{figure}[H]
\centering
\includegraphics[width=0.8\textwidth,height=0.8\textheight,keepaspectratio]{\fbpath/761.3371 КСВ/img/PROT_EMU_Q.pdf}
\captionof{figure}{Функциональная схема защиты ЭМУ от длительного протекания тока}\label{pic_ksv:PROT_EMU_Q}
\end{figure}

\item
Для электромагнитов управления предусмотрена защита от длительного протекания тока. Контроль
протекания тока через ЭМВ, ЭМО1 и ЭМО2 осуществляется при помощи токовых реле, положения контактов
которых должны быть заведены на соответствующие входные сигналы ФСУ \textit{«Работа ЭМВ»}, \textit{«Работа ЭМО1»},
\textit{«Работа ЭМО2»}. При длительном протекании тока по цепям ЭМО1 и ЭМВ формируется сигнал \textit{«Обест. ЭМО1
и ЭМВ»}, действующий на независимый расцепитель автомата питания цепей ЭМО1 и ЭМВ.

\item
Аналогично при длительном протекании тока по цепи ЭМО2 формируется сигнал \textit{«Обест. ЭМО2»} для
действия на независимый расцепитель автомата питания цепи ЭМО2. Регулируемые выдержки времени 
\textbf{«Тср защиты ЭМВ»}, \textbf{«Тср защиты ЭМО1»} и \textbf{«Тср защиты ЭМО2»} определяют время срабатывания на
обесточивание ЭМВ, ЭМО1 и ЭМО2 соответственно. Кроме того, наличие сигнала \textit{«Ср. Авар. давл. SF6»}
действует на обесточивание всех электромагнитов.

\begin{figure}[H]
\centering
\includegraphics[width=0.8\textwidth,height=0.8\textheight,keepaspectratio]{\fbpath/761.3371 КСВ/img/BLK_OP_Q.pdf}
\captionof{figure}{Функциональная схема блокировки управления выключателем}\label{pic_ksv:BLK_OP_Q}
\end{figure}

\item
На блокировку отключения выключателя действует сигнал \textit{«Ср. Авар. давл. SF6 В»}.

\item
На блокировку включения действуют сигналы \textit{«Ср. Авар. давл. SF6 В»}, \textit{«Неиспр. ЦУ В»}, \textit{«Ср. Пруж. не
заведены В»}, \textit{«Ср. Автомат ШП В»}.

\item
Также предусмотрена блокировка отключения от внешнего сигнала \textit{«Внеш. блок. откл. В»}, блокировка
включения от внешнего сигнала \textit{«Внеш. блок. вкл. В»} и блокировка управления (как включения, так и
отключения) от внешнего сигнала \textit{«Внеш. блок. упр. В»}.

\begin{figure}[H]
\centering
\includegraphics[width=0.8\textwidth,height=0.8\textheight,keepaspectratio]{\fbpath/761.3371 КСВ/img/Relay_fix_Q.pdf}
\captionof{figure}{Функциональная схема фиксации команд или включенного положения выключателя}\label{pic_ksv:Relay_fix_Q}
\end{figure}

\item
В устройстве предусмотрена программная фиксация команд или включенного положения выключателя
в зависимости от положения программного переключателя \textbf{«Реле фиксации В»}. В положении \textit{«РФК»}
фиксируется команда включения, в положении \textit{«РФП»} – включенное положение. Квитирование сигнала
фиксации \textit{«РФ В»} может осуществляться от сброса сигнализации или подачей оперативной команды на
отключение. При отключении выключателя от защит и наличии сигнала \textit{«РФ В1»} формируется сигнал о факте
аварийного отключения выключателя \textit{«Сигнал несоотв. В»}, необходимый для работы АПВ.

\small
\begin{longtable}{|>{\centering\arraybackslash}m{5.3cm}|>{\centering\arraybackslash}m{4.2cm}|>{\centering\arraybackslash}m{1cm}|>{\centering\arraybackslash}m{3.3cm}|>{\centering\arraybackslash}m{1.8cm}|}
% Первая страница: полный заголовок
\caption{Перечень уставочных параметров КСВ\hfill\vspace{-0.5\baselineskip}}\label{ksv:tbl1}\\
\hline
\rowcolor{gray!30}
Наименование уставки & Диапазон & Шаг & Значение по умолчанию & Единица измерения \\
\hline
\endfirsthead

% Промежуточные страницы: «Продолжение таблицы»
\caption*{\hspace{3pt}\emph{Продолжение таблицы \ref{ksv:tbl1}}\hfill\ }\vspace{-0.5\baselineskip} \\
\hline
\rowcolor{gray!30}
Наименование уставки & Диапазон & Шаг & Значение по умолчанию & Единица измерения \\
\hline
\endhead

%===m>КСВ|-
%===m
\end{longtable}
\normalsize	

\end{enumerate}